\documentclass{article}

\usepackage[utf8]{inputenc}
\usepackage[margin=1in]{geometry}
\usepackage{amsmath}
\usepackage{booktabs}
\usepackage{longtable}
\usepackage{graphicx}
\usepackage{xcolor}
\usepackage{hyperref}
\usepackage{listings}
\usepackage{float}
\usepackage{enumitem}
\usepackage{array}

\graphicspath{{../outputs/full_uci/}{outputs/full_uci/}{./}}

\hypersetup{
  colorlinks=true,
  linkcolor=blue,
  urlcolor=blue,
  citecolor=blue
}

\definecolor{codebg}{RGB}{248,248,248}
\definecolor{codekeyword}{RGB}{0,0,150}
\definecolor{codestring}{RGB}{150,70,20}
\definecolor{codecomment}{RGB}{80,120,80}

\lstdefinestyle{pythonstyle}{
  language=Python,
  backgroundcolor=\color{codebg},
  basicstyle=\ttfamily\small,
  keywordstyle=\color{codekeyword}\bfseries,
  stringstyle=\color{codestring},
  commentstyle=\color{codecomment},
  showstringspaces=false,
  breaklines=true,
  frame=single,
  numbers=left,
  numberstyle=\tiny,
  tabsize=4
}

\lstdefinestyle{textstyle}{
  backgroundcolor=\color{codebg},
  basicstyle=\ttfamily\small,
  showstringspaces=false,
  breaklines=true,
  frame=single,
  numbers=none,
  tabsize=2
}

\lstset{style=pythonstyle}

\newcommand{\safeimage}[2]{
\IfFileExists{#1}{
\begin{figure}[H]
\centering
\includegraphics[width=0.85\linewidth]{#1}
\caption{#2}
\end{figure}
}{
\begin{figure}[H]
\centering
\fbox{\parbox{0.82\linewidth}{Image not found in this Overleaf project. Upload the corresponding generated PNG image to display this figure.}}
\caption{#2}
\end{figure}
}
}

\title{\textbf{Federated Learning Optimization}\\
\large A Communication-Aware and Fairness-Aware Federated Learning Pipeline on UCI HAR}
\author{Author Name Placeholder}
\date{\today}

\begin{document}

\maketitle

\begin{abstract}
This report explains a federated learning optimization project implemented in Python and PyTorch. The project trains a human activity recognition model using the UCI Human Activity Recognition Using Smartphones dataset. Each human subject is treated as one federated client. Instead of gathering all client data into one central training table, the project simulates federated learning: a global model is sent to selected clients, each client trains locally, and the server averages the resulting model weights. The project studies four optimization ideas: FedAvg convergence, communication cost, CVaR-style fairness for weak clients, and hyperparameter search using grid search and differential evolution. The final pipeline produces training summaries, fairness summaries, communication shadow-price results, search results, plots, and a draft report.
\end{abstract}

\newpage
\tableofcontents
\newpage

\section{Executive Summary}

\subsection{Problem Solved}
The project asks a practical machine learning question: can a federated learning system be accurate, communication-efficient, and fair to weak clients at the same time?

Federated learning is useful when raw data should not be moved into one central server. In this project, each person in the UCI HAR dataset is treated as one client. The model travels to the clients, trains locally, and only model weights return to the server.

\subsection{Dataset}
The project uses the UCI Human Activity Recognition Using Smartphones dataset. It contains smartphone accelerometer and gyroscope features from 30 volunteers. The task is to predict one of six activities:
\begin{enumerate}
  \item WALKING
  \item WALKING\_UPSTAIRS
  \item WALKING\_DOWNSTAIRS
  \item SITTING
  \item STANDING
  \item LAYING
\end{enumerate}

\subsection{Model and ML Approach}
The main model is a small multilayer perceptron called \texttt{HARMLP}. It maps 561 input features to 6 output activity classes:
\[
561 \rightarrow 128 \rightarrow 64 \rightarrow 6.
\]

The training algorithm is FedAvg, short for Federated Averaging. In every round, a subset of clients receives the global model, trains locally, and sends updated weights back. The server averages those weights to form the next global model.

\subsection{Final Outputs}
The cleaned final output folder is \texttt{outputs/full\_uci}. It contains:
\begin{itemize}
  \item \texttt{convergence\_mean.png}
  \item \texttt{cvar\_pareto.png}
  \item \texttt{shadow\_price.png}
  \item \texttt{ga\_vs\_grid.png}
  \item \texttt{fedavg\_summary.csv}
  \item \texttt{cvar\_summary.csv}
  \item \texttt{grid\_search.csv}
  \item \texttt{lp\_shadow\_price.csv}
  \item \texttt{lp\_shadow\_price.json}
  \item \texttt{ga\_result.json}
  \item \texttt{report\_draft.md}
\end{itemize}

\subsection{Main Results}
\begin{table}[H]
\centering
\begin{tabular}{ll}
\toprule
Metric & Value \\
\midrule
Mean FedAvg accuracy & 0.9339 \\
Mean FedAvg loss & 0.1700 \\
Mean worst-client accuracy & 0.6914 \\
Best CVaR alpha by worst-client accuracy & 0.9 or 0.95 \\
Best CVaR worst-client accuracy & 0.7037 \\
Best grid search policy & \(E=4, K=10, \eta=0.02\) \\
Best grid search fitness & 0.9082 \\
Differential evolution fitness & 0.3993 \\
Differential evolution evaluations & 36 \\
LP budgets with positive shadow price & 5 out of 10 \\
\bottomrule
\end{tabular}
\caption{Main results from the generated output files.}
\end{table}

Precision, recall, F1 score, confusion matrix, saved model checkpoints, and example predictions are not available from the provided files. To report them, the project would need to save predictions, labels, and trained model weights.

\section{Repository Structure}

\subsection{File Tree}
\begin{lstlisting}[style=textstyle]
federated-learning-optimization/
  README.md
  requirements.txt
  .gitignore

  flopt/
    __init__.py
    config.py
    data.py
    models.py
    fedavg.py
    duality.py
    search.py
    plots.py

  experiments/
    run_full.py

  data/
    UCI HAR Dataset/
      README.txt
      features.txt
      features_info.txt
      activity_labels.txt
      train/
        X_train.txt
        y_train.txt
        subject_train.txt
      test/
        X_test.txt
        y_test.txt
        subject_test.txt

  outputs/
    full_uci/
      convergence_mean.png
      cvar_pareto.png
      shadow_price.png
      ga_vs_grid.png
      fedavg_summary.csv
      cvar_summary.csv
      grid_search.csv
      lp_shadow_price.csv
      lp_shadow_price.json
      ga_result.json
      report_draft.md
\end{lstlisting}

\subsection{Folder and File Purpose}
\begin{longtable}{p{0.28\linewidth}p{0.35\linewidth}p{0.27\linewidth}}
\toprule
Path & What it contains & Why it exists \\
\midrule
\texttt{README.md} & Short project summary and run command. & Helps a new user understand the project quickly. \\
\texttt{requirements.txt} & Python dependencies. & Makes environment setup reproducible. \\
\texttt{.gitignore} & Ignore rules for generated folders. & Prevents committing virtual environments, data, and outputs. \\
\texttt{flopt/config.py} & Federated learning configuration class. & Keeps training settings organized. \\
\texttt{flopt/data.py} & UCI HAR loader. & Converts raw files into federated clients. \\
\texttt{flopt/models.py} & Model definitions. & Defines the neural networks. \\
\texttt{flopt/fedavg.py} & FedAvg training and evaluation. & Core training engine. \\
\texttt{flopt/duality.py} & LP and KKT analysis. & Computes communication shadow prices. \\
\texttt{flopt/search.py} & Grid search and differential evolution. & Finds good training settings. \\
\texttt{flopt/plots.py} & Plot functions. & Creates report figures. \\
\texttt{experiments/run\_full.py} & Main experiment script. & Runs the full pipeline. \\
\texttt{data/UCI HAR Dataset} & Dataset files. & Provides real activity data. \\
\texttt{outputs/full\_uci} & Final outputs. & Stores metrics, plots, and report draft. \\
\bottomrule
\end{longtable}

\section{End-to-End Pipeline Overview}

\subsection{Workflow}
\begin{lstlisting}[style=textstyle]
Raw UCI HAR text files
    -> NumPy arrays
    -> Standardized features
    -> 30 ClientData objects
    -> HARMLP model
    -> FedAvg training
    -> Evaluation metrics
    -> CVaR fairness sweep
    -> Grid search and differential evolution
    -> Communication-budget LP
    -> Plots, CSV files, JSON files, and report draft
\end{lstlisting}

\subsection{Stage-by-Stage Explanation}
\begin{longtable}{p{0.24\linewidth}p{0.38\linewidth}p{0.28\linewidth}}
\toprule
Stage & What happens & Code file \\
\midrule
Initialization & The script parses command-line arguments and creates the output folder. & \texttt{experiments/run\_full.py} \\
Data loading & UCI HAR text files are read into arrays. & \texttt{flopt/data.py} \\
Preprocessing & Features are standardized using training-set statistics. & \texttt{flopt/data.py} \\
Client creation & Each subject becomes one federated client. & \texttt{flopt/data.py} \\
Model definition & The HAR MLP is created. & \texttt{flopt/models.py} \\
Training & FedAvg trains the global model. & \texttt{flopt/fedavg.py} \\
Evaluation & Average loss, average accuracy, and worst-client accuracy are computed. & \texttt{flopt/fedavg.py} \\
Fairness & CVaR-style client weighting is tested. & \texttt{flopt/fedavg.py} \\
Search & Grid search and differential evolution search over \(E,K,\eta\). & \texttt{flopt/search.py} \\
Duality & LP solves communication-budget policy selection. & \texttt{flopt/duality.py} \\
Results & CSV, JSON, PNG, and report files are written. & \texttt{run\_full.py}, \texttt{plots.py} \\
\bottomrule
\end{longtable}

\section{Dataset and EDA}

\subsection{Dataset Source}
The UCI HAR dataset contains sensor recordings from 30 volunteers aged 19 to 48. Each volunteer wore a smartphone on the waist. Accelerometer and gyroscope signals were sampled at 50 Hz. The original dataset provides 561 engineered features from time and frequency domains.

\subsection{Available Files}
\begin{itemize}
  \item \texttt{train/X\_train.txt}: training features.
  \item \texttt{train/y\_train.txt}: training labels.
  \item \texttt{train/subject\_train.txt}: subject IDs for train rows.
  \item \texttt{test/X\_test.txt}: test features.
  \item \texttt{test/y\_test.txt}: test labels.
  \item \texttt{test/subject\_test.txt}: subject IDs for test rows.
  \item \texttt{activity\_labels.txt}: class name mapping.
\end{itemize}

\subsection{Dataset Shape}
\begin{table}[H]
\centering
\begin{tabular}{lr}
\toprule
Property & Value \\
\midrule
Training rows & 7352 \\
Test rows & 2947 \\
Feature columns & 561 \\
Activity classes & 6 \\
Train subjects & 21 \\
Test subjects & 9 \\
Total subjects & 30 \\
Missing values in train features & 0 \\
Missing values in test features & 0 \\
Duplicate train feature rows & 0 \\
Feature minimum before project scaling & -1.0 \\
Feature maximum before project scaling & 1.0 \\
\bottomrule
\end{tabular}
\caption{UCI HAR dataset facts computed from the local files.}
\end{table}

\subsection{Label Distribution}
\begin{table}[H]
\centering
\begin{tabular}{llrr}
\toprule
Class ID & Activity & Train Count & Test Count \\
\midrule
1 & WALKING & 1226 & 496 \\
2 & WALKING\_UPSTAIRS & 1073 & 471 \\
3 & WALKING\_DOWNSTAIRS & 986 & 420 \\
4 & SITTING & 1286 & 491 \\
5 & STANDING & 1374 & 532 \\
6 & LAYING & 1407 & 537 \\
\bottomrule
\end{tabular}
\caption{Activity class distribution.}
\end{table}

\subsection{EDA Interpretation}
The classes are not perfectly balanced, but there is no extreme imbalance. The features have no missing values in the loaded files. The dataset already contains engineered numeric features, so no text tokenization, image processing, or embedding generation is needed.

\subsection{Recommended EDA Code}
\begin{lstlisting}
import numpy as np
from pathlib import Path

root=Path("data/UCI HAR Dataset")
x_train=np.loadtxt(root/"train/X_train.txt")
y_train=np.loadtxt(root/"train/y_train.txt",dtype=int)
subjects=np.loadtxt(root/"train/subject_train.txt",dtype=int)

print("Shape:",x_train.shape)
print("Missing values:",np.isnan(x_train).sum())
print("Feature min:",x_train.min())
print("Feature max:",x_train.max())

labels,counts=np.unique(y_train,return_counts=True)
for label,count in zip(labels,counts):
    print(label,count)

subject_ids,subject_counts=np.unique(subjects,return_counts=True)
print("Number of subjects:",len(subject_ids))
print("Rows per subject:",dict(zip(subject_ids,subject_counts)))
\end{lstlisting}

This code reads the train data, checks shape, checks missing values, shows the feature range, counts labels, and counts how many rows each subject contributes.

\section{Preprocessing}

\subsection{Label Conversion}
The UCI labels are numbered 1 through 6. PyTorch cross-entropy expects class labels from 0 through 5. Therefore, the code subtracts 1:
\begin{lstlisting}
y_train=np.loadtxt(root/"train"/"y_train.txt",dtype="int64")-1
y_test=np.loadtxt(root/"test"/"y_test.txt",dtype="int64")-1
\end{lstlisting}

\subsection{Feature Scaling}
\begin{lstlisting}
scaler=StandardScaler().fit(x_train)
x_train=scaler.transform(x_train).astype("float32")
x_test=scaler.transform(x_test).astype("float32")
\end{lstlisting}

The scaler is fitted only on training features. This avoids using test statistics to transform train data. The transformed arrays are converted to \texttt{float32}, which is the common tensor type for PyTorch neural networks.

\section{Model and Algorithms}

\subsection{HARMLP Architecture}
\begin{lstlisting}
class HARMLP(nn.Module):
    def __init__(self,features:int=561,classes:int=6):
        super().__init__()
        self.net=nn.Sequential(
            nn.Linear(features,128),
            nn.ReLU(),
            nn.Linear(128,64),
            nn.ReLU(),
            nn.Linear(64,classes),
        )

    def forward(self,x:torch.Tensor)->torch.Tensor:
        return self.net(x)
\end{lstlisting}

The model receives one 561-dimensional feature vector and outputs six raw scores called logits. The predicted class is the index of the largest logit.

\subsection{FedAvg Objective}
The global federated objective is:
\[
F(w)=\sum_{k=1}^{K}\frac{n_k}{n}F_k(w),
\]
where \(K\) is the number of clients, \(n_k\) is the number of examples on client \(k\), \(n\) is the total number of examples, and \(F_k(w)\) is the loss on client \(k\).

\subsection{FedAvg Aggregation}
After local training, the server averages selected client weights:
\[
w_{t+1}=\sum_{k \in S_t}\frac{n_k}{\sum_{j \in S_t}n_j}w_k.
\]

\section{Training}

\subsection{Main Training Configuration}
\begin{table}[H]
\centering
\begin{tabular}{ll}
\toprule
Setting & Value \\
\midrule
Dataset & UCI HAR \\
Clients & 30 \\
Rounds & 30 \\
Seeds & 7, 11, 19 \\
Default local epochs & 2 \\
Default clients per round & 10 \\
Default learning rate & 0.03 \\
Batch size & 32 \\
Optimizer & SGD \\
Loss & Cross entropy \\
\bottomrule
\end{tabular}
\caption{Main experiment configuration.}
\end{table}

\subsection{Local Client Training}
\begin{lstlisting}
def train_one_client(model,client,cfg,device):
    model.train()
    x=torch.tensor(client.x_train,dtype=torch.float32)
    y=torch.tensor(client.y_train,dtype=torch.long)
    loader=DataLoader(TensorDataset(x,y),batch_size=cfg.batch_size,shuffle=True)
    opt=torch.optim.SGD(model.parameters(),lr=cfg.lr)
    loss_fn=nn.CrossEntropyLoss()
    for _ in range(cfg.local_epochs):
        for xb,yb in loader:
            xb=xb.to(device)
            yb=yb.to(device)
            opt.zero_grad()
            loss=loss_fn(model(xb),yb)
            loss.backward()
            opt.step()
\end{lstlisting}

This is ordinary neural network training, but it happens separately on each selected client. The model makes predictions, the loss measures mistakes, backpropagation computes gradients, and SGD updates weights.

\subsection{Loss Function}
For one example with logits \(z\) and true label \(y\), cross-entropy is:
\[
\mathcal{L}(z,y)=-\log\frac{\exp(z_y)}{\sum_{c=1}^{C}\exp(z_c)}.
\]
High loss means the model is wrong or uncertain. Low loss means the model gives a high score to the correct class.

\section{Evaluation Metrics}

\subsection{Accuracy}
\[
\text{Accuracy}=\frac{\text{number of correct predictions}}{\text{number of predictions}}.
\]

\subsection{Average Client Accuracy}
\[
\text{MeanAccuracy}=\frac{1}{K}\sum_{k=1}^{K}\text{Accuracy}_k.
\]
This treats each client equally.

\subsection{Worst-Client Accuracy}
\[
\text{WorstClientAccuracy}=\min_k \text{Accuracy}_k.
\]
This metric is important because high average accuracy can hide one or two badly performing clients.

\section{CVaR Fairness}

\subsection{Idea}
CVaR-style weighting gives extra importance to high-loss clients. The code computes a quantile threshold:
\[
\tau=\text{Quantile}_{\alpha}(\ell_1,\ell_2,\ldots,\ell_K).
\]
Then it computes the tail amount:
\[
\text{tail}_k=\max(\ell_k-\tau,0).
\]
Clients with larger tail values receive more aggregation weight.

\subsection{Code}
\begin{lstlisting}
def _aggregation_weights(sizes,losses,cfg):
    weights=sizes.astype("float64")/sizes.sum()
    if cfg.cvar_alpha<=0:
        return weights
    tau=np.quantile(losses,cfg.cvar_alpha)
    tail=np.maximum(losses-tau,0)
    if tail.sum()==0:
        return weights
    weights=weights*(1+cfg.fairness_strength*tail/tail.sum())
    return weights/weights.sum()
\end{lstlisting}

\subsection{CVaR Results}
\begin{table}[H]
\centering
\begin{tabular}{rrrrr}
\toprule
Alpha & Accuracy & Accuracy Std & Worst Accuracy & Worst Accuracy Std \\
\midrule
0.00 & 0.9352 & 0.0033 & 0.6996 & 0.0354 \\
0.50 & 0.9351 & 0.0039 & 0.6955 & 0.0407 \\
0.75 & 0.9356 & 0.0038 & 0.6996 & 0.0354 \\
0.90 & 0.9366 & 0.0035 & 0.7037 & 0.0363 \\
0.95 & 0.9366 & 0.0035 & 0.7037 & 0.0363 \\
\bottomrule
\end{tabular}
\caption{CVaR sweep results.}
\end{table}

\safeimage{cvar_pareto.png}{Average client accuracy versus worst-client accuracy for CVaR settings.}

\section{Communication-Budget LP}

\subsection{Communication Cost}
The code estimates communication with:
\[
\text{upload bytes}=\text{parameters}\times4\times\text{selected clients},
\]
\[
\text{download bytes}=\text{parameters}\times4\times\text{selected clients}.
\]
The factor 4 means each parameter is assumed to be a 32-bit float.

\subsection{LP Formulation}
The LP is over candidate training policies, not over neural network weights. If policy \(i\) has loss \(\ell_i\) and communication cost \(c_i\), then:
\[
\begin{aligned}
\min_x \quad & \sum_i \ell_i x_i \\
\text{s.t.}\quad & \sum_i x_i=1,\\
& \sum_i c_i x_i \leq B,\\
& x_i\geq0.
\end{aligned}
\]

The dual variable \(\lambda^\star\) attached to the budget constraint is interpreted as the shadow price of communication.

\subsection{KKT Checks}
The code checks primal feasibility, dual feasibility, complementary slackness, and a stationarity residual. Five out of ten budget levels had positive shadow price in the generated results.

\safeimage{shadow_price.png}{Communication budget versus LP dual variable.}

\section{Hyperparameter Search}

\subsection{Search Variables}
The project searches over:
\begin{itemize}
  \item \(E\): local epochs,
  \item \(K\): clients per round,
  \item \(\eta\): learning rate.
\end{itemize}

\subsection{Fitness Function}
\[
\text{fitness}=\text{loss}+\gamma\cdot\text{communication}+10\max(0,0.80-\text{accuracy}).
\]
Lower fitness is better.

\subsection{Grid Search Results}
\begin{table}[H]
\centering
\begin{tabular}{rrrrrr}
\toprule
\(E\) & \(K\) & LR & Accuracy & Loss & Fitness \\
\midrule
4 & 10 & 0.02 & 0.9045 & 0.2636 & 0.9082 \\
2 & 10 & 0.03 & 0.8836 & 0.3333 & 0.9779 \\
4 & 15 & 0.02 & 0.9041 & 0.2694 & 1.2363 \\
1 & 10 & 0.03 & 0.7749 & 0.6535 & 1.5493 \\
2 & 15 & 0.01 & 0.7613 & 0.8605 & 2.2148 \\
1 & 5 & 0.01 & 0.6546 & 1.1620 & 2.9388 \\
\bottomrule
\end{tabular}
\caption{Grid search results.}
\end{table}

\subsection{Differential Evolution Result}
The differential evolution output was:
\[
x=[4.8939,4.1545,0.0717],
\]
with fitness \(0.3993\) after 36 evaluations. This is approximately \(E=5\), \(K=4\), and learning rate \(0.0717\).

\safeimage{ga_vs_grid.png}{Best grid-search fitness versus differential evolution fitness.}

\section{Results}

\subsection{FedAvg Results}
\begin{table}[H]
\centering
\begin{tabular}{rrrrr}
\toprule
Seed & Accuracy & Loss & Total Communication & Worst-Client Accuracy \\
\midrule
7 & 0.9272 & 0.1819 & 193396800 & 0.6296 \\
11 & 0.9355 & 0.1674 & 193396800 & 0.7037 \\
19 & 0.9391 & 0.1606 & 193396800 & 0.7407 \\
Mean & 0.9339 & 0.1700 & 193396800 & 0.6914 \\
\bottomrule
\end{tabular}
\caption{FedAvg summary over three seeds.}
\end{table}

\safeimage{convergence_mean.png}{Mean FedAvg convergence across seeds.}

\subsection{Interpretation}
The model performs well on average, with approximately 93.4 percent mean client accuracy. The worst-client accuracy is lower, approximately 69.1 percent, which shows that average accuracy alone is not enough. CVaR-style weighting slightly improves the weakest-client result. The LP analysis shows that communication can be treated as a constrained resource with a measurable shadow price.

\section{File-by-File Explanation}

\subsection{\texttt{flopt/config.py}}
This file defines \texttt{FLConfig}, a small immutable configuration object.
\begin{lstlisting}
@dataclass(frozen=True)
class FLConfig:
    rounds:int=20
    local_epochs:int=2
    clients_per_round:int=10
    lr:float=0.01
    batch_size:int=32
    seed:int=7
    cvar_alpha:float=0.0
    fairness_strength:float=1.0
\end{lstlisting}
If this file is removed, training scripts lose the central object that stores rounds, learning rate, batch size, seed, and fairness settings.

\subsection{\texttt{flopt/data.py}}
This file loads UCI HAR data and creates client datasets. Its most important class is:
\begin{lstlisting}
@dataclass
class ClientData:
    x_train:np.ndarray
    y_train:np.ndarray
    x_test:np.ndarray
    y_test:np.ndarray
\end{lstlisting}
Each \texttt{ClientData} object represents one subject. If this file is removed, the project cannot load data or build clients.

\subsection{\texttt{flopt/models.py}}
This file defines \texttt{HARMLP}, \texttt{LogisticModel}, \texttt{count\_parameters}, and \texttt{clone\_model}. The main executed model is \texttt{HARMLP}. If this file is removed, no trainable model exists.

\subsection{\texttt{flopt/fedavg.py}}
This is the training engine. It contains local training, global aggregation, evaluation, and CVaR-style weighting. If this file is removed, the project cannot train.

\subsection{\texttt{flopt/duality.py}}
This file solves the communication-budget LP using CVXPY. It also computes KKT diagnostics. If this file is removed, the shadow-price contribution disappears.

\subsection{\texttt{flopt/search.py}}
This file implements grid search and differential evolution. If this file is removed, the project cannot automatically compare training configurations.

\subsection{\texttt{flopt/plots.py}}
This file creates plots and configures Matplotlib to use local cache folders. If it is removed, the project can still train, but report-ready figures will not be created.

\subsection{\texttt{experiments/run\_full.py}}
This is the main entry point. It calls the loader, trainer, fairness sweep, search code, LP code, plot code, and report writer. If it is removed, the modules still exist, but there is no single command to run the complete project.

\section{How Files Call Each Other}
\begin{lstlisting}[style=textstyle]
experiments/run_full.py
  -> flopt/config.py
  -> flopt/data.py
  -> flopt/models.py
  -> flopt/fedavg.py
  -> flopt/search.py
  -> flopt/duality.py
  -> flopt/plots.py

flopt/search.py
  -> flopt/fedavg.py
  -> flopt/models.py

flopt/fedavg.py
  -> flopt/config.py
  -> flopt/data.py
  -> flopt/models.py
\end{lstlisting}

\section{Beginner Explanation}
Imagine 30 people each have phone sensor data. We want one model that recognizes activities, but we do not want to put everyone's data in one giant table. So the server sends a model to some people. Each person trains the model on their own data. Then they send back only the updated model weights. The server averages those weights. After many rounds, the shared model becomes better.

The model sees 561 numbers per example. Those numbers describe motion patterns. The model tries to learn which motion pattern means walking, sitting, standing, and so on.

The fairness part asks whether the model is helping the weakest clients. The communication part asks how expensive it is to send model weights. The search part asks which training settings give the best tradeoff.

\section{Issues, Risks, and Limitations}
\begin{itemize}
  \item Only three seeds were used.
  \item Raw per-round JSON logs were removed during cleanup.
  \item No saved model checkpoint is available.
  \item No precision, recall, F1 score, confusion matrix, or example predictions are available from the provided files.
  \item Differential evolution used a short evaluation budget.
  \item CVaR weighting is a practical heuristic, not a formal proof of CVaR convergence.
  \item Communication cost is estimated from parameter count, not measured network traffic.
\end{itemize}

\section{Suggested Improvements}
\begin{itemize}
  \item Add a dedicated EDA notebook.
  \item Save raw per-round logs for auditability.
  \item Save trained model checkpoints.
  \item Add confusion matrix, precision, recall, and F1 score.
  \item Add per-client metric tables.
  \item Run more seeds.
  \item Validate the best GA configuration using the full 30 or 100 rounds.
  \item Compare against FedProx or SCAFFOLD.
\end{itemize}

\section{Reproducibility Instructions}

\subsection{Install Dependencies}
\begin{lstlisting}[style=textstyle]
python -m venv .venv
.venv/bin/pip install -r requirements.txt
\end{lstlisting}

\subsection{Run Main Experiment}
\begin{lstlisting}[style=textstyle]
.venv/bin/python experiments/run_full.py --source uci --rounds 30 --seeds 7,11,19 --out outputs/full_uci --skip-ga
\end{lstlisting}

\subsection{Run With Differential Evolution}
\begin{lstlisting}[style=textstyle]
.venv/bin/python experiments/run_full.py --source uci --rounds 30 --seeds 7,11,19 --out outputs/full_uci
\end{lstlisting}

\subsection{Common Errors}
\begin{longtable}{p{0.32\linewidth}p{0.58\linewidth}}
\toprule
Error & Fix \\
\midrule
Dataset download fails & Manually place the UCI HAR dataset under \texttt{data/UCI HAR Dataset}. \\
Module not found & Run from the project root and install dependencies in \texttt{.venv}. \\
Matplotlib cache error & The project sets local cache paths in \texttt{plots.py}; delete cache folders and rerun if needed. \\
CVXPY solver error & Ensure CVXPY and its solvers are installed from \texttt{requirements.txt}. \\
Long runtime & Use fewer rounds or pass \texttt{--skip-ga}. \\
\bottomrule
\end{longtable}

\section{Appendix: Glossary}
\begin{longtable}{p{0.25\linewidth}p{0.65\linewidth}}
\toprule
Term & Meaning \\
\midrule
Federated learning & Training where data stays on clients and only model updates move. \\
Client & A local data holder. Here, one subject is one client. \\
Server & The coordinator that averages client model updates. \\
FedAvg & Federated Averaging, the main algorithm. \\
Round & One cycle of client selection, local training, and aggregation. \\
Local epoch & One pass over a client's local training data. \\
Loss & A number measuring how wrong the model is. \\
Cross entropy & Standard loss for multiclass classification. \\
Accuracy & Fraction of correct predictions. \\
Worst-client accuracy & The accuracy of the weakest client. \\
CVaR & A risk idea used here to emphasize high-loss clients. \\
LP & Linear program. \\
Dual variable & A number showing how important a constraint is. \\
Shadow price & The marginal value of relaxing a constraint. \\
KKT conditions & Mathematical checks for constrained optimization. \\
Grid search & Trying a fixed list of hyperparameters. \\
Differential evolution & A global search algorithm. \\
\bottomrule
\end{longtable}

\section{Appendix: Missing Items}
The following are not available from the provided files:
\begin{itemize}
  \item Saved final model weights.
  \item Example predictions.
  \item Confusion matrix.
  \item Precision, recall, and F1 score.
  \item Full retained per-round JSON logs after cleanup.
  \item Formal convergence theorem.
\end{itemize}

\end{document}
